\chapter{Background and Related Work}
\label{ch:background}

This chapter situates \voname{} in its technical and product context. It first
explains the concept of the \emph{virtual office} and why a spatial metaphor is
well suited to remote collaboration. It then reviews the enabling technologies
that make such a system possible --- real-time transport, peer-to-peer media,
game engines, and microservices --- so that the design decisions in later
chapters can be read against a shared vocabulary. Finally, it compares \voname{}
with existing tools and identifies the gap it aims to fill.

\section{The Virtual Office Concept}

A \emph{virtual office} (sometimes called a spatial collaboration platform or a
``metaverse office'') is a persistent, shared, two- or three-dimensional
environment in which remote team members are represented by avatars. Rather than
joining a meeting, a user \emph{enters a space}. Presence and communication are
governed by \emph{spatial proximity}: two users who move their avatars close to
one another can talk, exactly as they would in a physical room, and their
conversation is private to the people nearby.

This model differs fundamentally from conventional remote-collaboration tools in
three ways:

\begin{description}
  \item[Continuous presence, not discrete sessions.] In a video-conferencing
  tool, you are either in a call or not. In a virtual office, you are simply
  \emph{present in the space} for as long as you keep it open, and others can see
  where you are and what you are doing. Communication is a continuous ambient
  possibility rather than a scheduled event.

  \item[Proximity as the routing rule.] Who hears whom is determined by
  position. Walking toward a colleague opens a channel; walking away closes it.
  This reproduces the natural way conversations form and dissolve in a physical
  workplace, and it makes group boundaries (a meeting room, a huddle, a quiet
  corner) visible and intuitive.

  \item[Spatial memory and identity.] People associate activities with places ---
  ``the meeting room,'' ``my desk,'' ``the whiteboard by the window.'' A
  persistent map gives a distributed team the same spatial anchors, which aids
  orientation, belonging, and informal coordination.
\end{description}

The trade-off is engineering complexity. A virtual office is, at its heart, a
small \emph{multiplayer game}: it must synchronise the positions of many moving
avatars in real time, manage shared interactive objects, and start and stop live
media according to where those avatars are --- all while remaining responsive.
The remainder of this book is, in large part, an account of how \voname{} meets
that complexity.

\section{Enabling Technologies}

\subsection{Real-time Transport: WebSockets and STOMP}

Ordinary web traffic follows a request--response pattern over HTTP: the client
asks, the server answers, and the connection is (logically) done. This is a poor
fit for collaboration, where the server must \emph{push} events --- a new
message, a colleague moving, a notification --- to the client without being
asked. \textbf{WebSockets} solve this by upgrading an HTTP connection into a
persistent, bidirectional channel over which either side can send frames at any
time.

On top of raw WebSockets, \voname{} uses \textbf{STOMP} (Simple Text Oriented
Messaging Protocol), a lightweight publish/subscribe protocol. STOMP gives the
raw socket a structure of \emph{destinations}: clients \code{SUBSCRIBE} to
topics such as \route{/topic/channel/42} and \code{SEND} messages to application
destinations such as \route{/app/chat/send}, while the server broadcasts to
subscribers or to a specific user's private queue. This maps cleanly onto the
domain --- a channel is a topic, a user's error stream is a personal queue ---
and it is supported natively by Spring on the server and by client libraries in
JavaScript (\code{@stomp/stompjs}) and Dart (\code{stomp\_dart\_client}), both of
which \voname{} uses.

\begin{infobox}[title=The WebSocket Authentication Problem]
A browser cannot attach custom headers (such as an \code{Authorization} bearer
token) to a WebSocket handshake. \voname{} solves this uniformly across its
real-time services with a \emph{ticket} pattern: the client makes an
authenticated REST call to obtain a short-lived, single-use ticket, then opens
the socket with the ticket as a query parameter. This recurring pattern is
described in detail in Chapter~\ref{ch:architecture}.
\end{infobox}

\subsection{Peer-to-peer Media: WebRTC, PeerJS, and Agora}

Carrying live audio and video is expensive and latency-sensitive. Routing every
stream through the application's own servers would be costly and would not scale.
The industry standard for browser media is \textbf{WebRTC}, which establishes
direct, encrypted, peer-to-peer connections between clients for audio, video,
and arbitrary data, using the application server only for the initial
\emph{signalling} (helping peers discover and describe each other).

\voname{} uses two complementary media stacks:

\begin{itemize}
  \item In the 2D office, proximity calls are built on \textbf{PeerJS}, a thin
  wrapper over WebRTC that assigns each peer an id and simplifies the
  call/answer handshake. The office reuses the multiplayer session id as the
  PeerJS id, so one identity threads both systems.

  \item In the web client's meeting rooms, media is handled by \textbf{Agora}, a
  commercial real-time-communication platform whose Selective Forwarding Unit
  (SFU) relays streams efficiently. Here the application server acts as a
  \emph{token server}, minting short-lived, privilege-scoped access tokens so
  that only authorised users can join a room's media channel --- while the media
  itself never touches the application's servers.
\end{itemize}

The distinction matters: peer-to-peer (PeerJS/WebRTC) is simple and free but
scales poorly beyond small groups, whereas an SFU (Agora) scales to larger rooms
at the cost of a third-party dependency. \voname{} deliberately uses each where
it fits.

\subsection{Game Engines and Authoritative Servers}

The animated office is a real game, and it is built with a real game engine.
\textbf{Phaser} is a mature 2D HTML5 game framework providing sprite rendering,
tile-map support, arcade physics (for collisions and overlap detection),
animations, and camera control --- everything needed to render a walkable,
collision-aware office in the browser.

Multiplayer state is managed by \textbf{Colyseus}, an authoritative multiplayer
server framework. ``Authoritative'' means the server holds the single source of
truth for the shared world state; clients send their intentions and receive the
resulting state. Colyseus's schema system serialises only the \emph{changed
fields} of the state as compact binary diffs, so moving one avatar does not
require re-sending the entire world. This is the key to keeping many simultaneous
players responsive, and Chapter~\ref{ch:skyoffice} examines the resulting
synchronisation model, including client-side prediction and remote interpolation,
in detail.

\subsection{Microservices Architecture}

Rather than a single large application (a \emph{monolith}), \voname{}'s backend
is decomposed into independent \textbf{microservices}, each owning one business
capability (identity, chat, rooms, workspaces, and so on) and its own database.
Services communicate through well-defined interfaces --- synchronous REST for
queries that need an immediate answer, and asynchronous events over a message
broker for changes that other services should react to but need not wait for.

This style brings several benefits that recur through the book: each service can
use the persistence technology that best fits its data (a theme called
\emph{polyglot persistence}); services can be developed, deployed, and scaled
independently; and a failure in one service can be contained rather than taking
down the whole system. It also introduces challenges --- distributed data
consistency, cross-service authentication, and operational complexity --- whose
solutions form much of Chapter~\ref{ch:architecture}. The supporting cast of
technologies includes \textbf{Spring~Boot} and \textbf{Spring Cloud Gateway}
(the service and gateway framework), \textbf{RabbitMQ} (the message broker for
events), and a mix of \textbf{PostgreSQL}, \textbf{MongoDB}, and \textbf{Redis}
for storage.

\section{Related Work and Existing Tools}

\voname{} sits at the intersection of two product categories --- spatial/virtual
offices and unified team-collaboration suites --- and borrows from both. The
most relevant points of comparison are summarised below.

\begin{description}
  \item[Gather (gather.town).] The best-known commercial virtual office: a 2D,
  game-like map with avatars and proximity video, aimed at remote teams and
  events. Gather validated the spatial model at scale. \voname{} adopts the same
  core interaction (avatars, proximity media, interactive objects) but pairs it
  with a full, self-hosted collaboration backend rather than being purely a
  meeting space.

  \item[SkyOffice.] An open-source Phaser + Colyseus virtual-office project.
  \voname{}'s 2D world is an adapted fork of SkyOffice, which provides a strong,
  well-architected base for avatar movement, interactive computers and
  whiteboards, and proximity communication. The contribution of \voname{} is not
  the 2D engine in isolation but its integration into a broader, multi-client,
  microservices platform with real identity, workspaces, and persistent chat.

  \item[Slack and Microsoft Teams.] The dominant unified-collaboration suites.
  They excel at persistent, searchable messaging (channels, DMs, threads) and
  integrate calls, files, and apps --- but they are fundamentally
  \emph{list-based}, not spatial: presence is a coloured dot and communication is
  a discrete call or message. \voname{}'s chat subsystem is deliberately modelled
  on the Slack workspace/channel/thread structure (as later chapters show), but
  it is embedded in a spatial world that these tools lack.

  \item[Discord.] Combines persistent text channels with always-available voice
  ``rooms'' that users drop into and out of. \voname{}'s room service is
  explicitly Discord-like: a room bundles a persistent text channel with a
  joinable voice/video space. The difference again is the spatial layer and the
  workspace-centric organisation.

  \item[Zoom / Google Meet.] Best-in-class scheduled video conferencing, but
  session-based and non-persistent: there is no shared space, no ambient
  presence, and no spontaneous encounter. \voname{} treats real-time media as one
  capability among many rather than the whole product.
\end{description}

Table~\ref{tab:comparison} positions \voname{} against these tools along the
dimensions that matter for the problem stated in Chapter~\ref{ch:introduction}.

\begin{table}[htbp]
\centering
\caption{Comparison of \voname{} with representative existing tools.}
\label{tab:comparison}
\small
\begin{tabularx}{\textwidth}{L{3.0cm} C{1.7cm} C{1.7cm} C{1.9cm} C{1.7cm} L{2.4cm}}
\toprule
\textbf{Capability} & \textbf{Gather} & \textbf{Slack / Teams} & \textbf{Discord} & \textbf{Zoom} & \textbf{\voname{}} \\
\midrule
Spatial 2D world       & \checkmark & --- & --- & --- & \checkmark \\
Proximity media        & \checkmark & --- & partial & --- & \checkmark \\
Persistent channels    & partial & \checkmark & \checkmark & --- & \checkmark \\
Threads \& mentions    & --- & \checkmark & partial & --- & \checkmark \\
Meeting rooms          & \checkmark & \checkmark & \checkmark & \checkmark & \checkmark \\
Tasks \& calendar      & --- & \checkmark & --- & --- & planned \\
Self-hosted backend    & --- & --- & --- & --- & \checkmark \\
Native mobile client   & \checkmark & \checkmark & \checkmark & \checkmark & partial \\
\bottomrule
\end{tabularx}
\end{table}

\section{The Gap and This Project's Position}

Existing tools each solve part of the problem. Spatial offices deliver presence
and spontaneity but are typically thin on persistent, structured collaboration
and are closed, hosted products. Unified suites deliver excellent messaging and
integration but are list-based and lack a sense of place. Video tools deliver
media quality but nothing else.

\voname{} occupies the intersection: a spatial, game-like office \emph{and} a
full, self-hosted collaboration backend --- messaging, rooms, workspaces, and the
scaffolding for tasks and calendars --- surfaced through multiple clients. The
engineering value of the project lies in showing how these very different
concerns (a real-time multiplayer game, a Slack-style chat system, a
microservices identity and workspace layer, and peer-to-peer media) can be
integrated into one coherent platform. The chapters that follow describe that
integration from requirements through architecture to implementation.
